\exercisegroup

\begin{enumerate}
\def\labelenumi{\arabic{enumi})}
\item
  Let X be a topological space and let \(x, y\) be points of X. Prove that the maps \(p_x\) and \(p_y\), constant with respective images \(x\) and \(y\), are homotopic if and only if \(x\) and \(y\) belong to the same path-connected component of X.
\item
  Let S be the graph of the function from \([0,1]\) to R given by \(x \mapsto \sin(\pi/x)\) ; let \(\overline{S}\) be its closure in \(R^{2}\) .
\end{enumerate}

\begin{enumerate}
\def\labelenumi{\alph{enumi})}
\item
  Prove that the spaces S and \(\overline{S}\) are connected.
\item
  Prove that a connected and compact subset of \(\overline{S}\) containing the points \((1,0)\) and \((0,1)\) is necessarily equal to S, but that there does not exist a homeomorphism of I onto \(\overline{S}\) mapping 0 to \((1,0)\) and 1 to \((0,1)\) .
\item
  Determine the path-connected components of the space \(\overline{S}\) .
\item
  Prove that the space \(\overline{\mathrm{S}}\) is not locally connected.
\end{enumerate}

\begin{enumerate}
\def\labelenumi{\arabic{enumi})}
\setcounter{enumi}{2}
\item
  Let \(a \in \overline{\mathrm{S}}\) . Show that the pointed space \((\overline{\mathrm{S}}, a)\) has a universal covering if and only if \(a \in \{0\} \times [-1, 1]\) .
\item
  Let \(D = [0, 1] \cap (\mathbf{R} - \mathbf{Q})\) .
\end{enumerate}

\begin{enumerate}
\def\labelenumi{\alph{enumi})}
\item
  Let \(\varphi\colon D \to [0,1]\) be a map and let X be the complement in \([0,1]^2\) of the set formed by the points \((x,\varphi(x))\) for \(x\in\mathrm{D}\) . Show that X is connected.
\item
  Prove that there exists a bijective map \(\Phi\) from D onto the set of paths starting at \((0,0)\) and ending at \((1,1)\) in \([0,1]^2\) . For all \(x \in \mathrm{D}\) , let \(\psi(x)\) be a point with abscissa \(x\) on the path \(\Phi(x)\) . Let Y be the complement in \([0,1]^2\) of the set formed by the points \(\psi(x)\) for \(x \in \mathrm{D}\) . Show that the space Y is connected and locally connected, but that it is not path-connected.
\end{enumerate}

\begin{enumerate}
\def\labelenumi{\arabic{enumi})}
\setcounter{enumi}{4}
\tightlist
\item
  Let r be an element of \(N \cup \{\infty, \omega\}\) , let n and d be integers, and let X be a submanifold of class \(C^{r}\) of \(R^{n}\) , closed, purely of dimension \(d\). \(^{3}\)
\end{enumerate}

\begin{enumerate}
\def\labelenumi{\alph{enumi})}
\item
  Show that there exists a neighborhood U of X in \(\mathbf{R}^n\) and an isomorphism of class \(\mathrm{C}^r\) from \(\mathrm{X} \times ] - 1; 1[^{n - d}\) onto \(\mathrm{U}\) .
\item
  Every path in X is strictly homotopic to a path of class \(C^{r}\) .
\item
  If two paths of class \(C^{r}\) in X are homotopic, they are connected by a homotopy of class \(C^{r}\) .
\end{enumerate}

\begin{enumerate}
\def\labelenumi{\arabic{enumi})}
\setcounter{enumi}{5}
\item
  Let U be an open subset of a numerical space \(R^{n}\) . A path \(c: I \to U\) is said to be piecewise affine if there exists a finite sequence \((t_{0}, \ldots, t_{n})\) such that \(0 = t_{0} < t_{1} < \cdots < t_{n} = 1\) and such that the restriction of c to each interval \([t_{i-1}, t_{i}]\) is affine. Prove that every path in U is strictly homotopic to a piecewise affine path.
\item
  \begin{enumerate}
  \def\labelenumii{\alph{enumii})}
  \tightlist
  \item
    Let R be a set. If S is a subset of R, we denote by \(j_{S}\) the map from \(I^{S}\) into \(I^{R}\) that maps a function \(f: S \to I\) to the unique function from R to I whose restriction to S equals f and which is zero at every point of R-S.
  \end{enumerate}
\end{enumerate}

Prove that \(j_{S}\) is a homeomorphism from \(I^{S}\) onto the subspace of \(I^{R}\) consisting of functions vanishing outside S.

\begin{enumerate}
\def\labelenumi{\alph{enumi})}
\setcounter{enumi}{1}
\tightlist
\item
  Let X be the set of mappings from R to I that vanish outside a countable subset of R. We denote by Y the inductive limit of the family of spaces \(I^{S}\) , where S ranges over the set of countable subsets of R, and it is equipped with the inductive limit topology of the product topologies.
\end{enumerate}

The canonical maps from \(I^{S}\) into \(I^{R}\) induce a continuous bijection j from Y into X. This bijection is not a homeomorphism if R is not countable.

\begin{enumerate}
\def\labelenumi{\alph{enumi})}
\setcounter{enumi}{2}
\tightlist
\item
  Prove that, for every path \(c: I \to X\) there exists a countable subset S of R such that \(c(t)(s) = 0\) for all \(t \in I\) and every \(s \in R - S\) .
\end{enumerate}

\(d)\) Deduce that the map \(j^{-1}\) is continuous in the sense of arcs.

\begin{enumerate}
\def\labelenumi{\arabic{enumi})}
\setcounter{enumi}{7}
\tightlist
\item
  Let X be a connected compact topological space of cardinality \(\geqslant 2\) . We assume that for every pair \((x,y)\) of distinct points of X, the space \(X - \{x,y\}\) is not connected.
\end{enumerate}

\begin{enumerate}
\def\labelenumi{\alph{enumi})}
\item
  Show that for every \(x \in \mathrm{X}\) , the space \(\mathrm{X} - \{x\}\) is connected.
\item
  Let a, b be distinct points of X and let U, V be nonempty open and closed subsets of X such that \(X - \{a,b\}=U\cup V\) . Show that we have \(\overline{U} = U \cup \{a, b\}\) , \(\overline{V} = V \cup \{a, b\}\) and that these sets are connected.
\item
  We assume that X has a countable dense subset. Show then that X is homeomorphic to the circle \(S_{1}\) . (Prove that there exists a homeomorphism from I onto \(\overline{U}\) mapping 0 to a and 1 to b.)
\end{enumerate}

\begin{enumerate}
\def\labelenumi{\arabic{enumi})}
\setcounter{enumi}{8}
\tightlist
\item
  Let A be an uncountable well-ordered set every segment of which \([u, v]\) is countable.
\end{enumerate}

\begin{enumerate}
\def\labelenumi{\alph{enumi})}
\item
  Show that the set \(L = A \times [0, 1[\) , totally ordered by the lexicographic order and equipped with the topology \(\mathcal{T}_{0}(L)\) (TG, I, p. 91, exerc. 5), is connected (TG, IV, p. 48, exerc. 7). The space L is called the Alexandroff line.
\item
  Show that for every point \(x \in L\) , the topological space \(L - \{x\}\) is not connected.
\item
  Let \(\overline{L}\) be the ordered set obtained by adjoining to L points \(-\infty\) and \(+\infty\) such that \(-\infty < x < +\infty\) for all \(x \in L\) . Equip it with the topology \(\mathcal{T}_{0}(\overline{\mathrm{L}})\) ; it is called the completed Alexandroff line. Show that \(\overline{L}\) is compact and connected.
\end{enumerate}

d) Show that there does not exist a homeomorphism from \([0,1]\) onto \(\overline{\mathrm{L}}\) mapping 0 to \(-\infty\) and 1 to \(+\infty\) .

\begin{enumerate}
\def\labelenumi{\arabic{enumi})}
\setcounter{enumi}{9}
\tightlist
\item
  Let \(\overline{L}\) be the completed Alexandroff line and let S be the topological space obtained from \(\overline{L}\) by identifying the points \(-\infty\) and \(+\infty\) .
\end{enumerate}

\begin{enumerate}
\def\labelenumi{\alph{enumi})}
\item
  Show that the space S is compact and connected, of cardinal at least 2.
\item
  Show that for every pair \((x,y)\) of distinct points of S, the topological space \(S-\{x,y\}\) is not connected.
\item
  Show that the space S is not homeomorphic to a circle.
\end{enumerate}

\begin{enumerate}
\def\labelenumi{\arabic{enumi})}
\setcounter{enumi}{10}
\tightlist
\item
  \begin{enumerate}
  \def\labelenumii{\alph{enumii})}
  \tightlist
  \item
    Let X be the quotient space of the space \(\mathbf{R} \times \{1, 2\}\) by the finest equivalence relation for which \((t, 1)\) is equivalent to \((t, 2)\) if \(t \in \mathbf{R}^*\) . Show that there does not exist an injective path in X whose endpoints are the images of the points \((0, 1)\) and \((0, 2)\) .
  \end{enumerate}
\end{enumerate}

\begin{enumerate}
\def\labelenumi{\alph{enumi})}
\setcounter{enumi}{1}
\tightlist
\item
  Let \(c: I \to S^{1}\) be the path defined by \(c(t) = e^{3\pi it}\) ; it connects the point 1 to the point -1. Show that there is no injective path in \(S_{1}\) that is strictly homotopic to the path c.
\end{enumerate}

\begin{enumerate}
\def\labelenumi{\arabic{enumi})}
\setcounter{enumi}{11}
\item
  Let X be a paracompact space in which every point has a neighborhood homeomorphic to a complete metric space. Show that there exists a metric on X, compatible with its topology, that makes it a complete metric space (cf. TG, IX, p. 110, exerc. 34).
\item
  Let X be a separated topological space; we assume that every point of X has a neighborhood U such that every pair of points of U is joined by an injective path in U. If X is connected, show (without using proposition 18 of III, p. 282) that X is path-connected and that every pair of points of X is joined by an injective path.
\item
  Let X be a set and let c be an infinite cardinal such that c \textless{} Card(X).
\end{enumerate}

\begin{enumerate}
\def\labelenumi{\alph{enumi})}
\item
  Let \(\mathcal{U}_{\mathfrak{c}}\) be the set of subsets V of X such that \(\operatorname{Card}(X - V) \leqslant \mathfrak{c}\) or \(V = \varnothing\) . Prove that \(\mathcal{U}_{\mathfrak{c}}\) is a topology on \(X\) .
\item
  Show that the set X, equipped with this topology, is connected and locally connected, but that its path-connected components are reduced to a single element.
\item
  Show that the cone C(X) is path-connected, locally connected, but is not locally path-connected. This shows that in corollary 2 of III, p. 267, one cannot omit the hypothesis that the space is locally compact and metrizable.
\end{enumerate}

